% !TEX root = main.tex
\section{考前速查：结论与证明入口}
下面只列结论和起手步骤，完整论证见各题详解。证明题的“答案”包括理由，不能只写本表中的结论。
\small
\begin{longtable}{p{1.1cm}p{1.0cm}p{12.1cm}}
\toprule 章 & 题号 & 结论与起手步骤\\\midrule\endhead
零&4&依次：全成立；自反与传递；仅传递；对称与传递。\\
零&5&依次：仅对称；自反与对称；全不成立；仅对称。\\
零&6&递归定义与归纳；先证加法，再证分配律引理，最后证乘法。\\
零&7&对$c$归纳证明$a(b+c)=ab+ac$，右分配律由交换律推出。\\
零&8&若$x<y$，写$y=x+t$，$t>0$，则$ay=ax+at>ax$。\\
零&13&共同因子$d$使$a^{2^m}$同时同余$1$与$-1$，故$d\mid2$；取$a=2$构造互素数列。\\
零&14&无平方因子的约数对应不同素因子集合的子集，约数和为$(1-1)^r$。\\
零&23&$x\equiv17\pmod{31}$，逆元$27^{-1}\equiv23\pmod{31}$。\\
零&24&孙子定理原构造：$e_1=231,e_2=330,e_3=210$，$x=2e_1+e_2+5e_3\equiv302\pmod{385}$。\\\midrule
一&1&展开$(ab)^2=a^2b^2$，左右消去得到$ba=ab$。\\
一&2&所有元素等于自身逆元；$ab=(ab)^{-1}=b^{-1}a^{-1}=ba$。\\
一&4&解$ae=a$与$fa=a$；用任意元素的可解表示推广成双侧单位元，再构造逆元。\\
一&6&构造正$n$边形的二面体群$D_n$，阶$2n$，$n>2$时不交换。\\
一&7&共轭保持乘积，$a^{-1}b^ma=(a^{-1}ba)^m$，再对次数归纳。\\
一&8&取逆映射总是双射；保乘法当且仅当群交换。\\
一&9&等价关系当且仅当$S$为子群；等价类为右陪集$Sb$。\\
一&10&$\Z\to n\Z$，$k\mapsto nk$，是加法群同构。\\
一&12&指数$2$时，群外元素的左右陪集都等于$G\setminus H$。\\
一&18&无限阶元素会给出无穷多个不同循环子群；否则群是有限多个有限循环子群的并。\\
一&21&对满射$H\times K\to HK$计数，每个纤维与$H\cap K$一一对应。\\
一&22&正规性给出$MN=NM$；$m\mapsto mN$的核为$M\cap N$。\\
一&29&用$f(a)=1-a$把新运算转为旧运算；新零元$1$，新单位元$0$。\\
一&30&四组例子：$2\Z\subset\Z$；$\{0\}\times\Z\subset2\Z\times\Z$；角块子环；数量矩阵子环。\\
一&32&非零理想含$a\ne0$，再含$a^{-1}a=1$，故为全域。\\
一&34&若$r^n,s^m\in I$，展开$(r-s)^{n+m-1}$，每一项均落在$I$。\\
一&39&令$I=0$使用上题；非交换反例为$E_{12},E_{21}$。\\
一&43&真理想公共零点：紧致性加平方和；零因子：零点集有非空相对开区间。\\
一&44&$|M_n(\F_p)|=p^{n^2}$；$|\GL_n(\F_p)|=\prod_{k=0}^{n-1}(p^n-p^k)$。\\\midrule
二&7&有限个有理数统一分母后，问题化为整数的最大公因数。\\
二&8&有限阶生成元在交换群中能合并同类项，指数取有限范围。\\
二&10&对换生成$S_n$；用共轭把星形对换或相邻对换变成任意对换。\\
二&11&用长轮换共轭$(12)$，得到所有相邻对换。\\
二&12&依教材原题，奇数$n$用$(123),(12\cdots n)$，偶数$n$用$(123),(23\cdots n)$；共轭后逐点加入。\\
二&15&$(H_1\cap H_2)/N=(H_1/N)\cap(H_2/N)$，比较陪集代表元。\\
二&29&固定$g$，$A$与$gB^{-1}$的大小之和超过$|G|$，交集非空，故$g\in AB$。\\
二&30&按不同共轭子群计数，所有共轭共享单位元，覆盖数严格小于$|G|$。\\
二&31&$S_4$有3个西罗2子群和4个西罗3子群；分别列出8阶与3阶子群。\\
二&32&56阶群：若西罗7子群有8个，占48个非单位元；148阶群：西罗37子群唯一。\\
二&36&单位上三角群$U$阶为$p^{n(n-1)/2}$；正规化子为可逆上三角群$B$；个数为$\prod_{i=1}^n(p^i-1)/(p-1)^n$。\\
二&38&共轭同态$G\to\Aut(C_p)$的像阶同时整除$p$幂与$p-1$，故平凡。\\
二&39&$G/Z(G)$循环时，每个元素写成$a^iz$，两个元素直接相乘验证交换。\\
二&40&凯莱嵌入$S_n$，再以符号同态补上两个新点，嵌入$A_{n+2}$。\\
二&41&$[g,n]\in N\cap G^{(1)}=\{e\}$，故$n$与所有$g$交换。\\
二&42&子群的导群逐项包含；商群导群为原导群的像。\\
二&43&先找正规西罗子群；正规部分与商群均可解，使用可解性的扩张引理。\\
二&44&有限$p$群中心非平凡；商去中心后按群阶归纳。\\
二&45&映射$g\mapsto(gH,gK)$的核为$H\cap K$；商群嵌入两个可解群的直积。\\
二&46&$\Inn(G)\cong G/Z(G)$；若内自同构恰2个，则中心商循环，反而推出只有1个。\\
二&47&非交换群取非平凡内自同构；交换群先看取逆映射，指数2时交换两个基向量。\\
二&49&中心平凡；$S_3$自同构作用于三个对换，$S_4$自同构作用于四个西罗3子群的正规化子。\\\midrule
三&1&$(1-ba)^{-1}=1+b(1-ab)^{-1}a$，左右乘积均须验算。\\
三&2&固定$uv=1$，两个右逆之差满足$ut=0$；右逆唯一推出$vu=1$。\\
三&3&令$e=1-vu$，构造$x_k=v+eu^k$；若重复，右乘$v^j$推出$e=0$。\\
三&8&六项判断：对、错、错、错、对、错。反例须检验良定义、满射、保幺与非零条件。\\
\bottomrule
\end{longtable}
\normalsize
